% Generated draft: requires mathematical and source review.
\documentclass[11pt]{article}
\usepackage[margin=25mm]{geometry}
\usepackage{fontspec}
\setmainfont{Latin Modern Roman}
\usepackage{amsmath,amssymb,mathtools,unicode-math}
\setmathfont{Latin Modern Math}
\usepackage{enumitem,needspace,xurl,hyperref}
\setlength{\parindent}{0pt}
\setlength{\parskip}{5pt}
\setlength{\emergencystretch}{3em}
\newcommand{\N}{\mathbb{N}}
\newcommand{\Z}{\mathbb{Z}}
\newcommand{\Q}{\mathbb{Q}}
\newcommand{\R}{\mathbb{R}}
\newcommand{\C}{\mathbb{C}}
\newcommand{\F}{\mathbb{F}}
\newcommand{\A}{\mathbb{A}}
\newcommand{\PP}{\mathbb{P}}
\newcommand{\E}{\mathbb{E}}
\newcommand{\eps}{\varepsilon}
\newcommand{\dd}{\,\mathrm{d}}
\begin{document}
\section*{UnsolvedMath ID: 30001281}
\textbf{Source problem code:} catalogue entry
\subsection{Definitions and mathematical statement}
The supplied record describes \(F_\infty/F\) as a totally real \(p\)-adic Lie extension, with \(p\) a prime. Under the intended Galois interpretation, write \(G=\operatorname{Gal}(F_\infty/F)\). The usual interpretation of the Iwasawa algebra notation is
\[
\Lambda(G)=\Z_p[[G]]=\varprojlim_U\Z_p[G/U],
\]
where \(U\) ranges over the open normal subgroups of \(G\). The record does not specify the full hypotheses on the fields or extension, or confirm the coefficient ring used in this notation.

The notation \(\Lambda(G)_S\) denotes the intended localization at a multiplicative set \(S\), and \(K_1\) denotes algebraic \(K\)-theory in degree one. Neither \(S\) nor the hypotheses ensuring this localization are supplied. The symbol \(C^\bullet\) denotes the complex appearing in the source conjecture, and \(\partial\) its intended localization boundary map; the complex, its class group, and the domain and codomain of the boundary map are not fully defined in the record.

An Artin representation means a finite-dimensional continuous representation with finite image over a characteristic-zero coefficient field. The intended representations here are of \(G\), but their coefficient fields and the evaluation conventions are not supplied. The symbol \(\kappa_F\) is conventionally the \(p\)-adic cyclotomic character associated with \(F\). Its use as a character of \(G\), including the required factorization, is not specified. The expression \(\rho\kappa_F^k\) denotes the corresponding character twist when defined. The notation \(L_\Sigma(\rho,s)\) denotes the intended Artin \(L\)-function with the Euler factors indexed by \(\Sigma\) omitted; the set \(\Sigma\) and the conventions for these values and their comparison with \(p\)-adic evaluations are not supplied.

With these unresolved data understood as placeholders for the source conventions, the intended assertion is that there exists exactly one element
\[
\zeta\in K_1(\Lambda(G)_S)
\]
such that
\[
\partial(\zeta)=-[C^\bullet]
\]
and, for every intended Artin representation \(\rho\) and every integer \(k>0\) with \((p-1)\mid k\),
\[
\zeta(\rho\kappa_F^k)=L_\Sigma(\rho,1-k).
\]
The requested task is to construct this unique element. The record states uniqueness in the displayed \(K_1\)-group itself and supplies no quotient qualification.

This is an incomplete formulation pending recovery of the source hypotheses and definitions. Moreover, the supplied literature assessment describes the extraction as a remark about the abelian case, whereas the statement and its verification note describe a noncommutative conjecture. The intended restriction to abelian extensions, or inclusion of nonabelian extensions, therefore remains unresolved.
\subsection{Short English statement}
Construct the source's proposed unique \(p\)-adic zeta element satisfying both its boundary identity and its interpolation identity for every intended Artin representation and every positive integer divisible by \(p-1\). Essential hypotheses and definitions are missing, and the supplied descriptions disagree about the abelian versus nonabelian scope. The supplied open label is dated curation, not an independently verified status.
\subsection{Sources}
\begin{itemize}
\item Supplied local UnsolvedMath registry record, permanent ID \(30001281\), exact source problem code \texttt{catalogue entry}, title: Noncommutative \(p\)-adic Zeta Functions. The record belongs to the Oberwolfach Reports Open Problems collection and gives proposed year \(2009\); no individual proposer is supplied.
\item Supplied report citation: \emph{Algebraische Zahlentheorie} (2009), Oberwolfach Reports, pp.\ 1667--1730, DOI \texttt{10.4171/owr/2009/30}. Source URL and supplied literature source labeled ``OWR report'': \url{https://doi.org/10.4171/owr/2009/30}. The record describes these reports as documenting mathematical workshops and their open-problem discussions.
\item The registry labels the statement \texttt{corrected\_verified} and reports that it was checked against Conjecture 2.1 in the official report, retaining the boundary and interpolation conditions for the unique noncommutative zeta element. This is a supplied verification claim; no independent source consultation is asserted here.
\item Literature assessment dated exactly \texttt{2026-08-21}: the supplied status is ``open,'' and the supplied assessment says this is an extracted remark about the abelian case of a main conjecture, not a complete standalone resolution of the surrounding nonabelian problem. This assessment is supplied dated curation derived from the OWR material, not independently verified. The record expressly requests independent verification of time-sensitive status and cited claims.
\end{itemize}
\end{document}
